%%% source
%%% book: zorich-1
%%% chapter: 1
%%% section: 1
%%% title: Логическая символика
%%% pages: 1-4
%%%

\chapter*{Глава I. Некоторые общематематические понятия и обозначения}

\section{Логическая символика}

\subsection{Связки и скобки}

Язык этой книги, как и большинства математических текстов, состоит из обычного языка и ряда специальных символов излагаемых теорий. Наряду с этими специальными символами, которые будут вводиться по мере надобности, мы используем распространённые символы математической логики $\neg$, $\land$, $\lor$, $\Rightarrow$, $\Leftrightarrow$ для обозначения соответственно отрицания <<не>> и связок <<и>>, <<или>>, <<влечёт>>, <<равносильно>>\footnotemark[1].
\footnotetext[1]{В логике вместо символа $\land$ чаще используется символ \&. Символ $\Rightarrow$ импликации логики чаще пишут в виде $\to$, а отношение равносильности~--- в виде $\longleftrightarrow$ или $\leftrightarrow$. Однако мы будем придерживаться указанной в тексте символики, чтобы не перегружать традиционный для анализа знак $\to$ предельного перехода.}

Возьмем, например, три представляющих и самостоятельный интерес высказывания:

$L$. <<Если обозначения удобны для открытий \ldots, то поразительным образом сокращается работа мысли>> (Г.\,Лейбниц\footnotemark[2]).

\footnotetext[2]{Г.\,В.\,Лейбниц (1646--1716)~--- выдающийся немецкий учёный, философ и математик, которому наряду с Ньютоном принадлежит честь открытия основ анализа бесконечно малых.}

$P$. <<Математика~--- это искусство  называть разные вещи одинаковыми именами>> (А.\,Пуанкаре\footnotemark[3]).

\footnotetext[3]{А.\,Пуанкаре (1854--1912)~--- французский математик, блестящий ум которого преобразовал многие разделы математики и достиг её фундаментальных приложений в математической физике.}

$G$. <<Великая книга природы написана языком математики>> (Г.\,Галилей\footnotemark[4]).

\footnotetext[4]{Г.\,Галилей (1564--1642)~--- итальянский учёный, крупнейший естествоиспытатель. Его труды легли в основу всех последующих физических представлений о пространстве и времени. Отец современной физической науки.}

Тогда в соответствии с указанными обозначениями:

Запись $L \Rightarrow P$ означает \emph{$L$ влечёт $P$}

$L \Leftrightarrow P$ \qquad \emph{$L$ равносильно $P$}

$((L \Rightarrow P) \land (\neg P)) \Rightarrow (\neg L)$ \qquad \emph{Если $P$ следует из $L$ и $P$ неверно, то $L$ неверно}

$\neg((L \Leftrightarrow G) \lor (P \Leftrightarrow G))$ \qquad \emph{$G$ не равносильно ни $L$, ни $P$}

Мы видим, что пользоваться только формальными обозначениями, избегая разговорного языка,~--- не всегда разумно.

Мы замечаем, кроме того, что в записи сложных высказываний, составленных из более простых, употребляются скобки, выполняющие ту же синтаксическую функцию, что и при записи алгебраических выражений. Как и в алгебре, для экономии скобок можно договориться о <<порядке действий>>. Условимся с этой целью о следующем порядке приоритета символов:
\[
\neg, \quad \land, \quad \lor, \quad \Rightarrow, \quad \Leftrightarrow.
\]

При таком соглашении выражение $\neg A \land B \lor C \Rightarrow D$ следует расшифровать как $(((\neg A) \land B) \lor C) \Rightarrow D$, а соотношение $A \lor B \Rightarrow C$~--- как $(A \lor B) \Rightarrow C$, но не как $A \lor (B \Rightarrow C)$.

Записи $A \Rightarrow B$, означающей, что $A$ влечёт $B$ или, что то же самое, $B$ следует из $A$, мы часто будем придавать другую словесную интерпретацию, говоря, что $B$ есть \emph{необходимый признак} или \emph{необходимое условие} $A$ и, в свою очередь, $A$~--- \emph{достаточное условие} или \emph{достаточный признак} $B$. Таким образом, соотношение $A \Leftrightarrow B$ можно прочитать любым из следующих способов:

\emph{$A$ необходимо и достаточно для $B$};

\emph{$A$ тогда и только тогда, когда $B$};

\emph{$A$, если и только если $B$};

\emph{$A$ равносильно $B$}.

Итак, запись $A \Leftrightarrow B$ означает, что $A$ влечёт $B$ и, одновременно, $B$ влечёт~$A$.

Употребление союза \emph{и} в выражении $A \land B$ пояснений не требует.

Следует, однако, обратить внимание на то, что в выражении $A \lor B$ союз \emph{или} неразделительный, т.\,е. высказывание $A \lor B$ считается верным, если истинно хотя бы одно из высказываний $A$, $B$. Например, пусть $x$~--- такое действительное число, что $x^2 - 3x + 2 = 0$. Тогда можно написать, что имеет место следующее соотношение:
\[
(x^2 - 3x + 2 = 0) \Leftrightarrow (x = 1) \lor (x = 2).
\]

\subsection{Замечания о доказательствах}

Типичное математическое утверждение имеет вид $A \Rightarrow B$, где $A$~--- посылка, а $B$~--- заключение. Доказательство такого утверждения состоит в построении цепочки $A \Rightarrow C_1 \Rightarrow \ldots \Rightarrow C_n \Rightarrow B$ следствий, каждый элемент которой либо считается аксиомой, либо является уже доказанным утверждением\footnotemark[1].

\footnotetext[1]{Запись $A \Rightarrow B \Rightarrow C$ будет употребляться как сокращение для $(A \Rightarrow B) \land (B \Rightarrow C)$.}

В доказательствах мы будем придерживаться классического правила вывода: если $A$ истинно и $A \Rightarrow B$, то $B$ тоже истинно.

При доказательстве от противного мы будем использовать также принцип исключённого третьего, в силу которого высказывание $A \lor \neg A$ ($A$ или не $A$) считается истинным независимо от конкретного содержания высказывания $A$. Следовательно, мы одновременно принимаем, что $\neg(\neg A) \Leftrightarrow A$, т.\,е. повторное отрицание равносильно исходному высказыванию.

\subsection{Некоторые специальные обозначения}

Для удобства читателя и сокращения текста начало и конец доказательства условимся отмечать знаками $\blacktriangleleft$ и $\blacktriangleright$ соответственно.

Условимся также, когда это будет удобно, вводить определения посредством специального символа $:=$ (равенство по определению), в котором двоеточие ставится со стороны определяемого объекта.

Например, запись
\[
\int f(x)\,dx := \lim_{\lambda(P) \to 0} \sigma(f;\, P,\, \xi)
\]
определяет левую часть посредством правой части, смысл которой предполагается известным.

Аналогично вводятся сокращённые обозначения для уже определённых выражений. Например, запись
\[
\sum_{i=1}^{n} f(\xi_i)\,\Delta x_i =: \sigma(f;\, P,\, \xi)
\]
вводит обозначение $\sigma(f;\, P,\, \xi)$ для стоящей слева суммы специального вида.

\subsection{Заключительные замечания}

Отметим, что мы здесь говорили, по существу, только об обозначениях, не анализируя формализм логических выводов и не касаясь глубоких вопросов истинности, доказуемости, выводимости, составляющих предмет исследования математической логики.

Как же строить математический анализ, если мы не имеем формализации логики? Некоторое утешение тут может состоять в том, что мы всегда знаем или, лучше сказать, умеем больше, чем способны в данный момент формализовать. Пояснением смысла последней фразы может служить известная притча о сороконожке, которая даже ходить разучилась, когда её попросили объяснить, как именно она управляется со всеми своими конечностями.

Опыт всех наук убеждает нас в том, что считавшееся ясным или простым и нерасчленяемым вчера может подвергнуться пересмотру или уточнению сегодня. Так было (и, без сомнения, ещё будет) и с многими понятиями математического анализа, важнейшие теоремы и аппарат которого были открыты ещё в XVII--XVIII веках, но приобрели современный формализованный, однозначно трактуемый и, вероятно, потому общедоступный вид лишь после создания теории пределов и необходимой для неё логически полноценной теории действительных чисел (XIX век).

Именно с этого уровня теории действительных чисел мы и начнём в главе~II построение всего здания анализа.

Как уже отмечалось в предисловии, желающие быстрее ознакомиться с основными понятиями и эффективным аппаратом собственно дифференциального и интегрального исчисления могут начать сразу с главы~III, возвращаясь к отдельным местам первых двух глав лишь по мере необходимости.

\subsection*{Упражнения}

Будем отмечать истинные высказывания символом 1, а ложные~--- символом 0. Тогда каждому из высказываний $\neg A$, $A \land B$, $A \lor B$, $A \Rightarrow B$ можно сопоставить так называемую \emph{таблицу истинности}, которая указывает его истинность в зависимости от истинности высказываний $A$, $B$. Эти таблицы являются формальным определением логических операций $\neg$, $\land$, $\lor$, $\Rightarrow$. Вот они:

\begin{center}
\begin{tabular}{c|c}
 & $\neg A$ \\ \hline
$A=0$ & 1 \\
$A=1$ & 0
\end{tabular}
\qquad
\begin{tabular}{c|cc}
$A \land B$ & $B=0$ & $B=1$ \\ \hline
$A=0$ & 0 & 0 \\
$A=1$ & 0 & 1
\end{tabular}
\qquad
\begin{tabular}{c|cc}
$A \lor B$ & $B=0$ & $B=1$ \\ \hline
$A=0$ & 0 & 0 \\
$A=1$ & 0 & 1
\end{tabular}
\qquad
\begin{tabular}{c|cc}
$A \Rightarrow B$ & $B=0$ & $B=1$ \\ \hline
$A=0$ & 1 & 1 \\
$A=1$ & 0 & 1
\end{tabular}
\end{center}

\textbf{1.} Проверьте, все ли в этих таблицах согласуется с вашим представлением о соответствующей логической операции. (Обратите, в частности, внимание на то, что если $A$ ложно, то импликация $A \Rightarrow B$ всегда истинна.)

\textbf{2.} Покажите, что справедливы следующие простые, но очень важные и широко используемые в математических рассуждениях соотношения:

a) $\neg(A \land B) \Leftrightarrow \neg A \lor \neg B$; \qquad b) $\neg(A \lor B) \Leftrightarrow \neg A \land \neg B$; \qquad c) $(A \Rightarrow B) \Leftrightarrow (\neg B \Rightarrow \neg A)$;

d) $(A \Rightarrow B) \Leftrightarrow \neg A \lor B$; \qquad e) $\neg(A \Rightarrow B) \Leftrightarrow A \land \neg B$.
